\subsubsection{Haskell}
\begin{examdetails}
    typically either short coding task or proof of program
\end{examdetails}

\paragraph{Programming}
The best tip here is to read the Haskell book, and to solve the exercises during the semester.

\subparagraph{Fold}
One of the most important functions to understand is \texttt{foldr} (and \texttt{foldl}).
If you have used \texttt{reduce} functions before, in e.g. JavaScript / TypeScript,
they are similar to a \texttt{map} combined with a \texttt{reduce}.

For instance, in TypeScript, the \texttt{reduce} function has the type signature
\mint{typescript}|array.reduce( ( accumulator: A, current: A, idx: number, array: A[] ) => A, initialValue: A )|

In the Haskell prelude, they are defined as follows:
\begin{code}{haskell}
    foldr :: (a -> b -> b) -> b -> [a] -> b
    foldr f z []     = z
    foldr f z (x:xs) = f x (foldr f z xs)

    foldl :: (a -> b -> a) -> a -> [b] -> a
    foldl f z []     = z
    foldl f z (x:xs) = foldl f (f v x) xs
\end{code}

\subparagraph{zipWith}
To combine a \texttt{map} and a \texttt{zip} function, use \texttt{zipWith}, type:
\mint{haskell}|zipWith :: (a -> b -> c) -> [a] -> [b] -> [c]|

\TODO Add more remarks

\paragraph{Proofs}
These proofs use structural induction, see Section~\ref{sec:induction-proofs} for that.
